\documentclass[../../main.tex]{subfiles}
\begin{document}

\section{Product stress test of the all-cut Carleson estimate (Route A)}
\label{sec:product-stress}

This section executes Program~\ref{prog:product-test}: it stress-tests the all-cut absorptive
Carleson estimate (Assumption~\ref{ass:all-cut-carleson}) on product measures, where
(i) localization preserves product structure, (ii) KLS is known
(Proposition~\ref{prop:products}), and (iii) the operator norm of $A_t$ genuinely reaches
$\log n$, so a proof cannot route through $\lmax$-control and a failure would force the
geometric route.

Three core tools are pointed at the product model. First, the per-direction Carleson estimate
(Corollary~\ref{cor:per-direction}) is a budget:
$\E\int_0^\infty s_t\abs{G_t\theta}^2\dd t\le\theta^TR_0\theta\le1$ for every fixed direction.
Second, the Stein identity $\calS_\nu(E)=s^2\norm K_\HS^2$ (Proposition~\ref{prop:stein-rep})
makes the quadratic-chaos machinery of Section~\ref{sec:qcts} directly applicable to the
Riccati source. Third, the covariance technology of Section~\ref{sec:covariance-tech} supplies
the cut-free comparison quantities. Notation is that of this report throughout: $\mu_t$ is
Eldan stochastic localization \cite{Eldan2013ThinShell,LeeVempala2024},
$p_t,s_t,\delta_t,G_t,K_t,r_t,S_t,D_t,R_t$ are the two-color quantities, $\tau$ is the coarse
balanced exit time, and $X_t=(\lmax(A_t)-1)_+$.

We also record the structural facts about products from Proposition~\ref{prop:products}: if
$\mu=\bigotimes_{i=1}^n\mu^{(i)}$ with isotropic one-dimensional log-concave factors, then
$\mu_t$ is a product pathwise (the tilt factorizes), $A_t=\diag(A_t^{(1)},\dots,A_t^{(n)})$
with each $A^{(i)}$ a one-dimensional variance process of drift $-(A^{(i)})^2$, and
$h_{\mu_t}\ge c\,\lmax(A_t)^{-1/2}$ pathwise, so that KLS for products is not in question; the
question is whether the \emph{program's estimate} holds there.

\subsection{Theorem A: coordinate budgets, and the refutation of the rank-one candidate}
\label{sec:budgets}

\begin{lemma}[Block support]
\label{lem:block}
Let $\nu$ be a product probability measure on $\R^n$ and let $E$ be measurable with respect to
the coordinates in $J\subset\{1,\dots,n\}$. Then the two-color quantities of the pair
$(\nu,E)$ satisfy: $\delta_i=0$ for $i\notin J$, and $G_{ij}=K_{ij}=0$ unless both
$i,j\in J$.
\end{lemma}

\begin{proof}
For $i\notin J$, the coordinate $x_i$ is independent of $\one_E$, jointly with each $x_j$.
Hence $\E[x_i\mid E]=\E[x_i]$, giving $\delta_i=0$; for $j\in J$,
$\Cov(x_i,x_j\mid E)=\E[x_i]\,\E[x_j\mid E]-\E[x_i]\,\E[x_j\mid E]=0$, using
$\E[x_ix_j\one_E]=\E[x_i]\E[x_j\one_E]$; and for $j\notin J$, $j\ne i$,
$\Cov(x_i,x_j\mid E)=0=\Cov(x_i,x_j)$, while for $j=i$ the conditional variance is unchanged.
The same holds for $F=E^c$, so all entries of $G=\Sigma^E-\Sigma^F$ outside the $J\times J$
block vanish, and likewise for $K=G+(q-p)\delta\delta^T$.
\end{proof}

\begin{theorem}[Coordinate-budget theorem]
\label{thm:budget}
Let $\mu=\bigotimes_i\mu^{(i)}$ be a product of isotropic one-dimensional log-concave
measures, and let $E$ be measurable with respect to $k$ coordinates, with
$p_0\in[2/5,3/5]$. Then:
\begin{enumerate}[label=(\roman*),leftmargin=2.2em]
\item \emph{(Total source budget)}\quad
$\displaystyle \E\int_0^\infty S_t\dd t\ \le\ \sum_{i\in J}(R_0)_{ii}\ \le\ k$;
\item \emph{(Bounded information rate)}\quad $\E\,r_{t\wedge\tau}\le1+k$ for all $t$;
\item \emph{(Boundary bound)}\quad there is a universal $c>0$ with
\[
  \mu^+(E)\ \ge\ \frac{c}{\sqrt{1+k}}\;\min(p_0,q_0).
\]
\end{enumerate}
\end{theorem}

\begin{proof}
(i) Since $\mu_t$ is a product pathwise and $E$ is $J$-measurable, Lemma~\ref{lem:block}
applies to $(\mu_t,E)$ at every time: $G_t$ is supported on the $J\times J$ block, so
\[
  S_t=s_t\norm{G_t}_\HS^2=\sum_{i\in J}s_t\abs{G_te_i}^2 .
\]
The per-direction Carleson estimate Corollary~\ref{cor:per-direction} gives, for each fixed $i$
and every $T$,
\[
  \E\int_0^T s_t\abs{G_te_i}^2\dd t\le e_i^TR_0e_i=(R_0)_{ii}\le1,
\]
using $R_0=A_0-B_0\preceq A_0=I$. Sum over $i\in J$ and let $T\to\infty$ by monotone convergence.

(ii) The scalar Riccati identity Theorem~\ref{thm:scalar-riccati} gives, after stopping and taking
expectations within the regularity convention,
$\E r_{t\wedge\tau}\le r_0+\E\int_0^{t\wedge\tau}S_s\dd s\le1+k$, using $D\ge0$ and
$r_0\le1$ (rank-one $B_0\preceq I$).

(iii) By \eqref{eq:qv-p}, $\dd[p]_t=s_tr_t\dd t$ with $s_t\le\tfrac14$, so
\[
  \E[p]_{T\wedge\tau}\le\frac14\int_0^T\E r_{t\wedge\tau}\dd t\le\frac{(1+k)T}4 .
\]
Doob's $L^2$ inequality for the exit of the nested coarse window gives
$\Prob(\tau\le T)\le C_1(1+k)T$ with a universal $C_1$; choose $T_k=1/(2C_1(1+k))$ so that
$\Prob(\tau>T_k)\ge\tfrac12$, on which event $\min(p_{T_k},q_{T_k})\ge\tfrac13$. The
$T_k$-uniform log-concavity of $\mu_{T_k}$ and the perimeter supermartingale then give, as in
Theorem~\ref{thm:centroid-implies-kls},
$\mu^+(E)\ge\E\mu_{T_k}^+(E)\ge c\sqrt{T_k}\cdot\tfrac13\cdot\tfrac12\ge c'/\sqrt{1+k}$.
\end{proof}

\begin{corollary}[Refutation of the rank-one dynamic two-tail candidate]
\label{cor:refutation}
Let $\mu$ be a product as above and let $E$ be \emph{any} balanced cut depending on a single
coordinate $i$ --- in particular, any two-tail cut adapted to the coordinate whose variance
inflates under localization. Then
\[
  \E\int_0^\infty S_t\dd t\ \le\ 1
  \qquad\text{and}\qquad
  \mu^+(E)\ \ge\ c\,\min(p_0,q_0),
\]
with universal $c$. Consequently no such cut can witness a failure of the consumption chain of
Assumption~\ref{ass:all-cut-carleson}: the conclusion that the program derives from the Carleson estimate holds for
these cuts unconditionally. Quantitatively, the source spends only a vanishing fraction of time
at large height: for every \emph{deterministic} level $L>0$,
\[
  \E\,\bigl|\{t\ge0:\ S_t\ge\tfrac12L^2\}\bigr|
  \ \le\ \frac2{L^2}\,\E\int_0^\infty S_t\dd t\ \le\ \frac2{L^2}.
\]
Thus a two-tail configuration of source height $S_t\asymp\Lambda^2$
(Proposition~\ref{prop:two-tail} with $\Lambda=A_t^{(i)}$), wherever along the path the
coordinate variance $A_t^{(i)}$ inflates to order $\Lambda$, \emph{is self-extinguishing}: a
spike to height of order $\Lambda^2$ is paid at that rate from a per-coordinate budget of total
size $1$, so the time spent above any fixed fraction of it is $O(\Lambda^{-2})$.
\end{corollary}

\begin{proof}
Immediate from Theorem~\ref{thm:budget} with $k=1$; the occupation bound is Markov's inequality
at the deterministic level,
$\E\int_0^\infty\one_{\{S_t\ge L^2/2\}}\dd t\le\frac2{L^2}\E\int_0^\infty S_t\dd t$, together with
the total source budget $\E\int_0^\infty S_t\dd t\le1$ from Theorem~\ref{thm:budget}(i).
\end{proof}

\begin{remark}[What the refutation does and does not say]
\label{rem:refutation-scope}
Corollary~\ref{cor:refutation} verifies, for single-coordinate cuts, the \emph{total stopped
source bound} --- which is what the Riccati--Gronwall consumption of Theorem~\ref{thm:carleson-implies-centroid}
actually uses --- not the literal interval form of Assumption~\ref{ass:all-cut-carleson}; the distinction is
immaterial for the program, since the consumption needs only
$\E r_{t\wedge\tau}$ bounded, which (ii) provides directly. The corollary closes the specific
attack proposed in the original formulation of the test (``a cut adapted to the dynamically
large coordinate''): the per-direction estimate, which entered the program after that
proposal, is exactly the tool that kills it. The original ``teeth'' of the test --- a proof of
Assumption~\ref{ass:all-cut-carleson} cannot route through $\lmax(A_t)$ --- is hereby complemented on the
counterexample side: \emph{a counterexample cannot route through a single inflated coordinate
either}. Both the proof and the refutation are forced into genuinely high-dimensional,
cut-specific territory; Theorem~\ref{thm:budget} localizes the entire remaining danger of the
product model in cuts of unbounded coordinate complexity.
\end{remark}

\subsection{Theorem C: the covariance reduction for arbitrary cuts of products}
\label{sec:covariance-reduction}

For cuts of unbounded complexity the block structure is unavailable, but the product structure
yields a pathwise quadratic-chaos bound at every time --- not only at $t=0$ --- which converts
the all-cut estimate into a cut-free covariance statement.

\begin{lemma}[Quadratic chaos for non-isotropic products]
\label{lem:product-qcts}
Let $\nu=\bigotimes_i\nu^{(i)}$ be a product of centered one-dimensional log-concave measures
with variances $\sigma_i^2$, $A=\diag(\sigma_i^2)$, $Y\sim\nu$. Then for every symmetric $M$,
\begin{equation}\label{eq:product-qcts}
  \Var\bigl(Y^TMY\bigr)\ \le\ C_*\,\norm{A^{1/2}MA^{1/2}}_\HS^2 ,
\end{equation}
with $C_*$ universal.
\end{lemma}

\begin{proof}
Expand $Y^TMY=\sum_iM_{ii}Y_i^2+\sum_{i\ne j}M_{ij}Y_iY_j$. By independence and centering, all
cross-covariances vanish: $\Cov(Y_i^2,Y_iY_j)=\E[Y_i^3]\E[Y_j]=0$,
$\Cov(Y_iY_j,Y_iY_k)=\E[Y_i^2]\E[Y_j]\E[Y_k]=0$ for $j\ne k$, and disjoint pairs are
independent. Hence, using $M_{ij}=M_{ji}$ so that each unordered pair $\{i,j\}$ contributes
$\Var\bigl(2M_{ij}Y_iY_j\bigr)=4M_{ij}^2\sigma_i^2\sigma_j^2$,
\[
  \Var(Y^TMY)=\sum_iM_{ii}^2\Var(Y_i^2)
  +\sum_{i<j}4M_{ij}^2\sigma_i^2\sigma_j^2
  =\sum_iM_{ii}^2\Var(Y_i^2)+2\sum_{i\ne j}M_{ij}^2\sigma_i^2\sigma_j^2 .
\]
For one-dimensional log-concave $Y_i$, the reverse H\"older inequalities (see e.g.\
\cite[Cor.~5]{KLnotes}) give $\E Y_i^4\le C_4\sigma_i^4$ with $C_4$ universal, so
$\Var(Y_i^2)\le(C_4-1)\sigma_i^4$. Both sums are dominated by
$C_*\sum_{ij}M_{ij}^2\sigma_i^2\sigma_j^2=C_*\norm{A^{1/2}MA^{1/2}}_\HS^2$ with
$C_*=\max(C_4-1,2)$.
\end{proof}

\begin{theorem}[Pathwise covariance bound on the source]
\label{thm:covariance-bound}
Let $\mu$ be a product as in Theorem~\textup{\ref{thm:budget}} and $E$ \emph{any} measurable
set. Then pathwise, on the coarse balanced window $\{p_t\in[1/3,2/3]\}$,
\begin{equation}\label{eq:S-le-X2}
  S_t\ \le\ C\bigl(1+X_t^2\bigr),
\end{equation}
with $C$ universal.
\end{theorem}

\begin{proof}
Fix $t$ and write $\nu=\mu_t$, centered at $a_t$, with diagonal covariance $A=A_t$; $\nu$ is a
product by Proposition~\ref{prop:products}. By the Stein representation
Proposition~\ref{prop:stein-rep}, $\Cov_\nu(\one_E,Y^TMY)=s\inner KM$ for $Y=x-a_t$. Cauchy--Schwarz
and Lemma~\ref{lem:product-qcts} give
\[
  s^2\inner KM^2\le\Var(\one_E)\,\Var(Y^TMY)\le s\,C_*\norm{A^{1/2}MA^{1/2}}_\HS^2 .
\]
Substituting $M=A^{-1/2}NA^{-1/2}$ and taking the supremum over $\norm N_\HS\le1$ yields the
whitened bound
\begin{equation}\label{eq:whitened}
  s\,\norm{A^{-1/2}KA^{-1/2}}_\HS^2\ \le\ C_* .
\end{equation}
Un-whitening, $\norm K_\HS\le\lmax(A)\,\norm{A^{-1/2}KA^{-1/2}}_\HS$, so
$s\norm K_\HS^2\le C_*\lmax^2$. Finally $G=K-(q-p)\delta\delta^T$ and, on the coarse window
($s\ge2/9$, $\abs{q-p}\le1$),
\[
  S=s\norm G_\HS^2\le2s\norm K_\HS^2+2s\abs\delta^4
  \le2C_*\lmax^2+\frac{2r^2}{s}\le2C_*\lmax^2+9\lmax^2 ,
\]
using $r\le\lmax$ (rank-one $B\preceq A$). Since $\lmax\le1+X$ and $(1+X)^2\le2(1+X^2)$, the
claim follows.
\end{proof}

\begin{corollary}[Covariance reduction of the all-cut estimate for products]
\label{cor:V2-implies}
Define the \emph{second-moment interface condition} on a window $[0,T_0]$:
\begin{equation}\label{eq:V2}
  \tag{V2}
  \int_I\E X_t^2\dd t\ \le\ K\,\abs I
  \qquad\text{for every interval }I\subset[0,T_0].
\end{equation}
If \textup{(V2)} holds for the localization of a product measure $\mu$, then
Assumption~\ref{ass:all-cut-carleson} holds for $\mu$ on $[0,T_0]$ with $C_0=C(1+K)$, $C_1=0$ and
damping coefficient $\alpha=0$. In particular the entire all-cut estimate for products reduces
to a statement about $n$ independent one-dimensional variance processes, with no reference to
cuts.
\end{corollary}

\begin{proof}
$\E\int_{I\cap[0,\tau]}S_t\dd t\le C\int_I(1+\E X_t^2)\dd t\le C(1+K)\abs I$ by
Theorem~\ref{thm:covariance-bound}, since the coarse window holds on $\{t<\tau\}$.
\end{proof}

\begin{remark}
This is the precise sense in which the product model isolates the covariance content of
Route~A: a cut-free second-moment bound suffices. Note the consistency with the discipline
warning of \eqref{eq:trivial-lambda}: (V2) concerns the \emph{excess} $X=(\lmax-1)_+$ in square,
not $\E\int\lmax\le CT$ for all measures (which would already imply KLS and is not available);
for products, (V2) is not circular --- KLS for products is known independently --- it is a
concrete, decidable question about explicit one-dimensional SDEs. The next section decides it
on the polylog window and assesses it beyond.
\end{remark}

\subsection{Theorem D: the second-moment window and the sharpness regime}
\label{sec:window}

\subsubsection{The second-moment bound on the polylog window}

The required covariance input is the small-time operator-norm control of
Section~\ref{sec:covariance-tech}: the sup-over-time window of Theorem~\ref{thm:KL-window} and
its discharge (Corollary~\ref{cor:KI-discharged}). It yields the second-moment interface bound
that \textup{(V2)} requires.

\begin{theorem}[Second-moment interface bound on the polylog window]
\label{thm:V2-window}
There are universal constants $c_0,K_0$ such that for every isotropic log-concave $\mu$ on
$\R^n$ (in particular every product) and every $t\le t_1(n):=c_0(\log n)^{-2}$,
\[
  \E\,X_t^2\ \le\ K_0 ,
  \qquad\text{hence \textup{(V2)} holds on }[0,t_1(n)]\text{ with }K=K_0 .
\]
Consequently, by Corollary~\textup{\ref{cor:V2-implies}}, Assumption~\ref{ass:all-cut-carleson} holds for product
measures on the window $[0,c_0(\log n)^{-2}]$, unconditionally.
\end{theorem}

\begin{proof}
On the event $\{\norm{A_t}_\op<2\}$, $X_t\le1$. On the complement, of probability at most
$e^{-1/(Ct)}$ by Theorem~\ref{thm:KL-window}, the Brascamp--Lieb cap
\eqref{eq:BL-cap} gives $X_t\le t^{-1}$ pathwise. Hence
\[
  \E X_t^2\le1+t^{-2}e^{-1/(Ct)}\le1+\sup_{u>0}u^2e^{-u/C}=1+(2C/e)^2=:K_0 .
\]
\end{proof}

\subsubsection{Sharpness, and the open window of the literature}

Two cited facts delimit what can be expected beyond $t_1(n)$. First, the window of
Theorem~\ref{thm:KL-window} is essentially sharp: this is asserted in
\cite[Remark~62]{KLnotes} and realized by the explicit example of \cite[Remark~64 and
Proposition~65]{KLnotes}: the product of two-sided exponentials, for
which, in Klartag's formulation, one of the eigenvalues ``can reach $\log n$ in a short
time'', so that the operator norm is at least of order $\log n$ at times of order
$1/\log n$ \cite{KlartagChenNotes}. Second --- and this is worth highlighting --- extending
the covariance bound to the intermediate window is explicitly open and explicitly valuable:
it is recorded in the literature around \cite{Guan2025Tail,KLnotes} that the validity of the
covariance bound for $t\in[(\log n)^{-2},(\log n)^{-1}]$ would already improve the best known
isoperimetric bound of \cite{Klartag2023Logarithmic}. The interface functionals
$\Xi_T$ and $\Xi^{(2)}_T$ of \eqref{eq:Xi2-def} therefore sit exactly at the recognized
frontier of the covariance technology, which is evidence that the difficulty has been
parametrized correctly rather than relocated.

\begin{heuristic}[Expected failure of (V2) on universal windows for exponential products]
\label{heur:V2-fails}
Let $\mu$ be the product of $n$ two-sided exponentials. Suppose, as the sharpness discussion
of \cite[Remark~62]{KLnotes} and \cite{KlartagChenNotes} indicates, that with probability
bounded below some coordinate variance reaches the scale of the Brascamp--Lieb cap,
$A_t^{(i)}\asymp1/t$, on a time window $[t_1,2t_1]$ with $t_1\asymp1/\log n$. Then on that
window $\E X_t^2\gtrsim t^{-2}$, and
\[
  \int_{t_1}^{2t_1}\E X_t^2\dd t\ \gtrsim\ \frac1{t_1}\ \asymp\ \log n .
\]
Hence on any \emph{universal} window $[0,T_0]$, the condition \textup{(V2)} is expected to
fail for exponential products, with $\int_0^{T_0}\E X_t^2\dd t\gtrsim\log n$. We mark this as
a heuristic: the cited statements establish that the threshold $2$ is crossed after the
polylog window and that the bound is ``pretty much sharp''; the quantitative strength
$\lmax\asymp1/t$ with non-negligible probability requires the sharpness example to be carried
out in detail, which we have not done line by line.
\end{heuristic}

\begin{remark}[Consequence of the heuristic, stated conditionally]
\label{rem:covariance-route-dead}
If Heuristic~\ref{heur:V2-fails} is correct, then the covariance-only route of
Corollary~\ref{cor:V2-implies} cannot prove Assumption~\ref{ass:all-cut-carleson} for products on a universal window
--- even though both the assumption's conclusion (KLS for products) and the per-cut estimates
of Theorem~\ref{thm:budget} are true there. This is the precise product-model analogue of the
crude-interface insufficiency of Remark~\ref{rem:insufficiency}: cut-free covariance information
saturates at a logarithm. The stress test then returns the verdict the original proposal
anticipated, in refined form: any proof of Assumption~\ref{ass:all-cut-carleson} for products must be
\emph{cut-aware}, and Theorem~\ref{thm:budget} exhibits the prototype mechanism --- fixed
per-coordinate budgets paid from the within-class variance $R_0$ --- while
Corollary~\ref{cor:refutation} shows that the dual counterexample mechanism must likewise be
cut-aware and high-complexity. The residue is a single question.
\end{remark}

\subsection{The residual: the adapted alignment problem}
\label{sec:residual}

Combining Theorems~\ref{thm:budget} and \ref{thm:covariance-bound}, the anatomy of any
counterexample to Assumption~\ref{ass:all-cut-carleson} within the product model is tightly constrained. To violate
$\E\int_0^{T\wedge\tau}S_t\dd t\le MT$ for a given large $M$ and universal $T$, a fixed
balanced cut $E$ must satisfy: (a) $E$ depends on at least $MT$ coordinates
(Theorem~\ref{thm:budget}(i)); (b) at least of order $MT$ units of per-coordinate budget
$\sum_i\E\int s\abs{Ge_i}^2$ --- each coordinate contributing at most $1$ over \emph{all}
time --- must be spent inside the common short window $[0,T]$; (c) the expenditure must be
concentrated, by Theorem~\ref{thm:covariance-bound}, at times when $X_t$ is large, i.e.\
aligned with the rare inflation excursions of the independent coordinate processes; and
(d) throughout, the mass must remain balanced ($t<\tau$) and the source must not be matched
by the damping $D_t$, which the Riccati identity subtracts for free. The cut is fixed before
the Brownian path; the inflating coordinate indices and the conditional means $a_{t,i}$
around which a two-tail configuration would have to center are random. The decisive question
is whether (b)--(d) are simultaneously achievable.

\begin{question}[Adapted alignment problem for products]
\label{q:alignment}
Let $\mu$ be a product of $n$ isotropic two-sided exponentials. Does there exist a universal
$T_0$ and constants $C_0,C_1,\alpha<1$ such that for every $n$ and every balanced cut $E$,
\[
  \E\int_{I\cap[0,\tau]}\ \sum_{i:\,A_t^{(i)}\ge2}s_t\abs{G_te_i}^2\dd t
  \ \le\ C_0\abs I+C_1\E\int_{I\cap[0,\tau]}r_t\dd t
  +\alpha\,\E\int_{I\cap[0,\tau]}D_t\dd t ,
  \qquad I\subset[0,T_0]\ ?
\]
Equivalently, up to the covariance-controlled part: can a fixed cut spend
$\Omega(1)$-fractions of unboundedly many per-coordinate budgets inside the inflation
excursions of those coordinates, within a common universal window, while staying balanced and
underdamped? An affirmative answer to the displayed estimate, combined with
Theorem~\textup{\ref{thm:V2-window}} for the polylog window and a matching treatment of the
non-inflated part at moderate times, would settle Program~\ref{prog:product-test} positively. A negative answer
--- an explicit high-complexity cut achieving the alignment --- would refute Assumption~\ref{ass:all-cut-carleson}
and force Route~B unconditionally.
\end{question}

\begin{remark}[Designated test family]
\label{rem:test-family}
The natural candidate cuts for the negative direction are permutation-symmetric
thin-shell-type sets $E=\{x:\sum_i\psi(x_i)\ge\theta\}$ with $\psi$ even, which have unbounded
coordinate complexity and, at $t=0$, bounded source by Lemma~\ref{lem:product-qcts}. The
question is whether localization can dynamically correlate their per-coordinate expenditures
with the inflation excursions. Symmetry is a double-edged constraint: it spreads the budget
uniformly ($\E\int_0^Ts\abs{Ge_i}^2$ independent of $i$), but the inflating index is
exchangeable too, so symmetric cuts are precisely the ones whose budget allocation a random
inflating coordinate ``finds'' automatically. A careful computation of
$\E\int_0^{T}S_t\dd t$ for this family --- one-dimensional log-concave estimates,
independence, and the explicit variance dynamics are all available --- is, in our assessment,
the single most informative computation now open on Route~A, and it is finite-dimensional in
difficulty in a way that the open targets of Section~\ref{sec:open} are not.
\end{remark}
\end{document}
